\documentclass[12pt]{article}

\usepackage{amsmath, amssymb, amsthm}
\usepackage{geometry}
\geometry{margin=1in}

\title{CSE400: Fundamentals of Probability in Computing\\
Lecture 6: Discrete Random Variables, Expectation and Problem Solving}
\author{Dhaval Patel, PhD\\
Associate Professor\\
Computer Science and Engineering (CSE)}
\date{January 22, 2025}

\begin{document}
\maketitle

\section*{Lecture Outline}
\begin{itemize}
    \item Random Variables (RVs): Definition and Concept
    \item Types of Random Variables
    \item Discrete Random Variables
    \item Probability Mass Function (PMF)
    \item Examples of Discrete Random Variables
    \item Bayes’ Theorem (Recap and Example)
\end{itemize}

\section{Random Variables}

\subsection{Definition}
A \textbf{random variable} $X$ on a sample space $\Omega$ is a function
\[
X : \Omega \rightarrow \mathbb{R}
\]
that assigns to each sample point $\omega \in \Omega$ a real number $X(\omega)$.

\subsection{Discrete Random Variables}
In this lecture, we restrict attention to \textbf{discrete random variables}.  
A random variable is said to be discrete if it takes values in a range that is finite or countably infinite.

Although $X$ is defined as a function mapping into $\mathbb{R}$, the actual set of values
\[
\{X(\omega) : \omega \in \Omega\}
\]
is a discrete subset of $\mathbb{R}$.

\section{Distribution of a Discrete Random Variable}

The distribution of a discrete random variable can be visualized using a bar diagram:
\begin{itemize}
    \item The $x$-axis represents the values that the random variable can take.
    \item The height of the bar at value $a$ represents $\Pr[X = a]$.
\end{itemize}

Each probability $\Pr[X = a]$ is computed by evaluating the probability of the corresponding event in the sample space.

\section{Types of Random Variables}

\subsection{Discrete Random Variables}
Characteristics:
\begin{itemize}
    \item Countable support
    \item Probability Mass Function (PMF)
    \item Probabilities are assigned to individual values
    \item Each possible value has strictly positive probability
\end{itemize}

\subsection{Continuous Random Variables}
Characteristics:
\begin{itemize}
    \item Uncountable support
    \item Probability Density Function (PDF)
    \item Probabilities are assigned to intervals
    \item Each individual value has probability zero
\end{itemize}

\section{Example 1: Tossing Three Fair Coins}

Suppose an experiment consists of tossing three fair coins.

Let $Y$ denote the number of heads that appear.

Then $Y$ is a random variable that can take values
\[
\{0, 1, 2, 3\}.
\]

\subsection*{Computation of Probabilities}

\[
\Pr(Y = 0) = \Pr((t,t,t)) = \frac{1}{8}
\]

\[
\Pr(Y = 1) = \Pr((t,t,h),(t,h,t),(h,t,t)) = \frac{3}{8}
\]

\[
\Pr(Y = 2) = \Pr((t,h,h),(h,t,h),(h,h,t)) = \frac{3}{8}
\]

\[
\Pr(Y = 3) = \Pr((h,h,h)) = \frac{1}{8}
\]

Since $Y$ must take one of the values $0$ through $3$, we have
\[
1 = \Pr\left(\bigcup_{i=0}^{3} \{Y = i\}\right)
  = \sum_{i=0}^{3} \Pr(Y = i).
\]

\section{Probability Mass Function (PMF)}

\subsection{Definition}
A random variable that can take at most a countable number of possible values is said to be \textbf{discrete}.

Let $X$ be a discrete random variable with range
\[
R_X = \{x_1, x_2, x_3, \dots\}
\]
where the range is finite or countably infinite.

The function
\[
P_X(x_k) = \Pr(X = x_k), \quad k = 1,2,3,\dots
\]
is called the \textbf{Probability Mass Function (PMF)} of $X$.

\subsection{Property of PMF}
Since $X$ must take one of the values $x_k$, we have
\[
\sum_{k=1}^{\infty} P_X(x_k) = 1.
\]

\section{PMF Example}

Let the probability mass function of a random variable $X$ be given by
\[
p(i) = c \frac{\lambda^i}{i!}, \quad i = 0,1,2,\dots
\]
where $\lambda > 0$.

\subsection*{Step 1: Find the constant $c$}

Since $\sum_{i=0}^{\infty} p(i) = 1$, we have
\[
c \sum_{i=0}^{\infty} \frac{\lambda^i}{i!} = 1.
\]

Using the identity
\[
e^{\lambda} = \sum_{i=0}^{\infty} \frac{\lambda^i}{i!},
\]
we obtain
\[
c e^{\lambda} = 1 \quad \Rightarrow \quad c = e^{-\lambda}.
\]

\subsection*{Step 2: Compute $\Pr(X = 0)$}
\[
\Pr(X = 0) = e^{-\lambda} \frac{\lambda^0}{0!} = e^{-\lambda}.
\]

\subsection*{Step 3: Compute $\Pr(X > 2)$}
\[
\Pr(X > 2) = 1 - \Pr(X \le 2)
\]
\[
= 1 - \Pr(X=0) - \Pr(X=1) - \Pr(X=2).
\]

\section{Bayes’ Theorem}

\subsection{Recall}
Using
\[
\Pr(A \cap B) = \Pr(B \mid A)\Pr(A),
\]
we obtain
\[
\Pr(B_i \mid A) = 
\frac{\Pr(A \mid B_i)\Pr(B_i)}
{\sum_{j=1}^{n} \Pr(A \mid B_j)\Pr(B_j)}.
\]

This is known as the \textbf{Bayes Formula (Proposition 3.1)}.

\subsection{Terminology}
\begin{itemize}
    \item $\Pr(B_i)$ is the \textbf{a priori probability}
    \item $\Pr(B_i \mid A)$ is the \textbf{posterior probability}
\end{itemize}

\section{Bayes’ Theorem Example: Auditorium with 30 Rows}

An auditorium has 30 rows of seats.

\begin{itemize}
    \item Row 1 has 11 seats
    \item Row 2 has 12 seats
    \item Row 3 has 13 seats
    \item $\dots$
    \item Row 30 has 40 seats
\end{itemize}

A door prize is awarded by:
\begin{enumerate}
    \item Randomly selecting a row (each row equally likely)
    \item Randomly selecting a seat within that row
\end{enumerate}

\subsection*{Given}
\[
\Pr(R_k) = \frac{1}{30}
\]

\subsection*{Compute}
\[
\Pr(S_{15} \mid R_{20}) = \frac{1}{30}
\]

\subsection*{Bayes’ Computation}
\[
\Pr(R_{20} \mid S_{15}) =
\frac{\Pr(S_{15} \mid R_{20}) \Pr(R_{20})}
{\sum_{k=1}^{30} \Pr(S_{15} \mid R_k)\Pr(R_k)}
\]

\end{document}
